\section{Methodology and System Design}
\label{sec:method}

\subsection{High-Level Architecture}

\begin{figure}[H]
  \centering
  \includegraphics[width=\linewidth]{fig_capstone_moe_framework.png}
  \caption{Mixture-of-experts schematic : input text, hierarchical reasoning block, parallel ML (TF--IDF + linear), DL (transformer encoders), and HRM experts, fused by a learned gating network.}
  \label{fig:capstone_moe}
\end{figure}

Four phases: (1)~preprocessing and splitting, (2)~expert training/inference, (3)~ensemble combination (stacking ensembles STACK1--7, gating MOE1--5), (4)~validation and reporting. The design favors heterogeneity to increase complementarity. Figure~\ref{fig:capstone_moe}  and Figure~\ref{fig:capstone_ensemble_pipe} (training/evaluation flow) illustrate the same decomposition for architectural clarity; formal expert IDs remain Table~\ref{tab:experts}.

\subsection{Preprocessing}
Deduplication and text cleaning (URLs, HTML, mentions, hashtags, contractions, case normalization); subword (BPE / WordPiece) tokenization per model family; stratified train/validation/test splits; $k$-fold cross-validation ($k=5$) to generate out-of-fold predictions for stacking without test leakage; binary and multi-class label remapping (e.g., coarse star ratings mapped to sentiment classes where applicable).

\subsection{Expert Families}
\begin{table*}[t]
  \centering
  \caption{Expert Families and Roles}
  \label{tab:experts}
  \begin{tabular}{@{}lll@{}}
    \toprule
    \textbf{Family} & \textbf{Role} & \textbf{Examples} \\
    \midrule
    ML & Fast baselines & B1, B2, E-ML1, E-ML2 \\
    DL & Sequential / CNN & BiLSTM+attn., CNN, E-DL$*$ \\
    Transformers & Contextual encoders & DistilBERT, BERT, RoBERTa (e.g., B3, B5, E-DL3) \\
    HRM & Multi-level reasoning & E-HRM1--E-HRM3 (levels L1--L4) \\
    \bottomrule
  \end{tabular}
\end{table*}

Note: ML baselines use TF-IDF (1-3 n-grams, capped vocabulary) with logistic regression or linear SVM and class weights where needed.

\subsection{HRM Formulation}
Lexical, syntactic, semantic, and pragmatic modules compose into a final representation:
$L_1 = \phi_1(x)$,
$L_2 = \phi_2(x, L_1)$,
$L_3 = \phi_3(x, L_1, L_2)$,
$L_4 = \phi_4(x, L_1, L_2, L_3)$,
$y_{\mathrm{HRM}} = f_{\mathrm{final}}(L_1, L_2, L_3, L_4)$.
In the thesis implementation, this hierarchy supports inspectable chains (sentiment-bearing terms, negation, discourse context, sarcasm or irony cues) even when raw HRM accuracy lags stronger experts.

\subsection{Ensemble Combination}
\textbf{Stacking:} train experts with 5-fold CV; collect OOF probabilities; train meta-learner (e.g., logistic regression); evaluate on held-out test without leakage.
\textbf{MoE gating:} dense softmax weights $g_i(x)$ over experts; $\hat{y} = \sum_i g_i(x)\, f_i(x)$; sparse top-$K$ variants for efficiency.
\textbf{Optimization:} weighted cross-entropy for class imbalance, AdamW with learning rates on the order of $2\times10^{-5}$ for transformer experts and $\sim\!10^{-3}$ for linear heads, warmup and decay schedules, early stopping on macro-F1, and gradient clipping where configured.

\subsection{Validation protocol}
Macro-F1 is the primary metric, supplemented by accuracy, precision, recall, and AUROC as reported in the thesis tables. Runs use multiple random seeds (e.g., 42, 123, 456) with paired tests, confidence intervals, and effect sizes where detailed statistical comparisons are required. Cross-domain checks (e.g., train on one corpus, evaluate on another) probe retention under shift, consistent with the original validation plan.

\subsection{Software Stack}
Python~3.9 or newer; PyTorch; Hugging Face Transformers; scikit-learn; Pandas. The thesis layout includes \path{config/}, \path{models/}, \path{train/}, \path{checkpoints/}, and evaluation scripts that write rollups under \path{output/metrics/}.

\subsection{Reporting scope}
The full thesis contains numbered algorithms for dataset exploration, the cleaning pipeline, stacking with OOF collection, MoE forward passes, HRM forward passes, and a generic training loop. This paper cites those artifacts through the implementation report rather than duplicating pseudocode blocks.
